\section{Introdução}
\label{sec_introducao}

A computação em nuvem~\citep{mell2011nist} transformou profundamente a forma como infraestruturas são provisionadas e gerenciadas. Provedores públicos como Amazon Web Services (AWS), Google Cloud Platform (GCP) e Microsoft Azure consolidaram um modelo operacional no qual máquinas virtuais e contêineres são criados, configurados e descartados de forma programática e automatizada. Central a esse modelo está o \textit{Instance Metadata Service} (IMDS), um serviço interno acessível por um endereço IP reservado, tipicamente \texttt{169.254.169.254}, pertencente à faixa de endereços \textit{link-local} IPv4 definida pela RFC~3927~\citep{rfc3927}. Por meio dele, cada instância descobre sua própria configuração em tempo de inicialização: nome do host, endereços IP, chaves SSH autorizadas e \textit{scripts} de \textit{user-data} \citep{aws_imds}. Essa capacidade é o alicerce sobre o qual ferramentas de provisionamento automatizado, especialmente o \textit{cloud-init}, operam para configurar instâncias sem intervenção manual \citep{cloudinit2024}.

Gerenciadores de contêineres e máquinas virtuais voltados à infraestrutura privada, como o Incus, projeto derivado (\textit{fork}) do LXD mantido pela organização Linux Containers \citep{incus2024, incus_announcement}, carecem de um serviço de metadados nativo equivalente ao oferecido pelas nuvens públicas. Essa ausência obriga administradores a configurar instâncias manualmente ou a recorrer a \textit{scripts} ad hoc que não seguem nenhum padrão estabelecido, rompendo o fluxo de provisionamento esperado pelo \textit{cloud-init} e inviabilizando a portabilidade de configurações entre ambientes públicos e privados. O resultado prático é que ambientes Incus, por mais capazes que sejam em termos de virtualização, tornam-se cidadãos de segunda classe no ecossistema de automação de infraestrutura.

Embora soluções como o serviço de metadados do OpenStack abordem esse problema no contexto de plataformas de nuvem privada completas \citep{openstack_metadata}, não existe uma solução leve e independente projetada especificamente para ambientes Incus. Isso é particularmente relevante para cenários de \textit{homelab}, computação de borda (\textit{edge computing})~\citep{shi2016} e nuvens privadas de pequena escala, nos quais implantar uma plataforma completa como o OpenStack representa uma complexidade operacional desproporcional às necessidades do ambiente. Existe, portanto, uma lacuna técnica não atendida pelas soluções existentes.

Este trabalho propõe o desenvolvimento de um serviço de metadados compatível com \textit{cloud-init} para o Incus. O serviço implementa uma API REST que entrega metadados, \textit{user-data}, \textit{vendor-data} e configuração de rede às instâncias, seguindo a interface esperada pelo \textit{cloud-init} \citep{cloudinit2024}. A arquitetura adota um modelo orientado a eventos combinado com sincronização periódica, garantindo que os dados das instâncias permaneçam atualizados de forma consistente com o estado reportado pelo gerenciador Incus. O serviço foi concebido para ser implantado de forma autônoma, sem dependência de componentes externos de grande porte, atendendo diretamente ao perfil de ambientes de infraestrutura privada de menor escala. Para cenários que exigem tolerância a falhas, o serviço oferece, opcionalmente, alta disponibilidade por meio do algoritmo de consenso RAFT~\citep{ongaro2014raft}, com múltiplos nós coordenados por eleição de líder e replicação de \textit{log}.

As principais contribuições deste trabalho são: (a) uma análise das implementações de serviços de metadados em nuvens públicas e privadas, identificando os requisitos mínimos para compatibilidade com o \textit{cloud-init}; (b) o projeto e a implementação de um serviço de metadados compatível com \textit{cloud-init} para o Incus, incluindo a modelagem da API, o mecanismo de sincronização de instâncias, a persistência dos dados e o suporte opcional a alta disponibilidade por consenso RAFT; e (c) a avaliação experimental do serviço em ambientes Incus reais, cobrindo a corretude do provisionamento automatizado, a latência e a escalabilidade da API com até 200 contêineres simultâneos e o comportamento do mecanismo de alta disponibilidade em um \textit{cluster} de três nós. O restante deste trabalho está organizado da seguinte forma: a Seção~\ref{sec_relacionados} apresenta os trabalhos relacionados; a Seção~\ref{sec_proposta} detalha a proposta e a implementação; a Seção~\ref{sec_resultados} apresenta a avaliação experimental; e a Seção~\ref{sec_conclusao} conclui o trabalho.
